\chapter{算术：因数与质数}\label{ch:g9:arith}

算术研究整数，以及整数之间如何相互整除。
其核心对象是质数，它们是通过乘法构成每个整数的基本积木。
本章以最大公因数收尾，
它是一劳永逸地化简分数的恰当工具。
这一故事还将在高中分册
及以后的学习中继续展开。

\section{因数与倍数}

\begin{definition}[因数与倍数]\label{def:g9:arith:divisor}
设 $a$ 与 $b$ 为正整数。称 $b$
\emph{整除}\index{整除} $a$（或称 $b$ 为 $a$ 的因数，或称
$a$ 是 $b$ 的\emph{倍数}\index{倍数}），若 $a = b \times
k$ 对某个整数 $k$ 成立 --- 即 $a$ 被 $b$ 除
所得余数为 $0$。
\end{definition}

\begin{example}
$24$ 的因数有 $1, 2, 3, 4, 6, 8, 12, 24$ --- 它们成对出现，
每一对之积都是 $24$：$(1,24)$、$(2,12)$、$(3,8)$、$(4,6)$。
$7$ 的倍数有 $7, 14, 21, 28, \dots$。
\end{example}

\begin{proposition}[整除判别法]\label{prop:g9:arith:rules}
一个整数在以下情形时能被相应的数整除：
\begin{itemize}
  \item 能被 $2$ 整除，当其末位数字为偶数（$0, 2, 4, 6, 8$）时；
  \item 能被 $5$ 整除，当其末位数字为 $0$ 或 $5$；
  \item 能被 $10$ 整除，当其末位数字为 $0$；
  \item 能被 $3$（或 $9$）整除，当其各位数字之和能被
        $3$（或 $9$）整除；
  \item 能被 $4$ 整除，当其末两位数字构成的数能被 $4$ 整除。
\end{itemize}
\end{proposition}
\admitted

\begin{example}
$7\,215$ 的末位是 $5$，故能被 $5$ 整除。其各位数字之和为
$7 + 2 + 1 + 5 = 15$，能被 $3$ 整除而不能被 $9$ 整除，故 $7\,215$
能被 $3$ 整除而不能被 $9$ 整除。事实上 $7\,215 = 3 \times 5 \times 481$。
\end{example}

\section{质数}

\begin{definition}[质数]\label{def:g9:arith:prime}
称\emph{质数}\index{质数}为一个 $\geq 2$ 的整数，它的
因数只有 $1$ 与它自身。小于 $30$ 的质数为
\[
2,\ 3,\ 5,\ 7,\ 11,\ 13,\ 17,\ 19,\ 23,\ 29 .
\]
按约定，数 $1$ \emph{不}是质数。
一个 $\geq 2$ 而不是质数的整数，称为\emph{合数}。
\end{definition}

\begin{theorem}[质因数分解]\label{thm:g9:arith:factorization}
每个 $\geq 2$ 的整数都可以写成质数的乘积，且
这种分解在不计因子次序的意义下是唯一的。
\end{theorem}
\admitted

\begin{method}[整数的质因数分解]\label{met:g9:arith:factorization}
反复用尽可能小的质数去除，直到得到 $1$：
\begin{enumerate}
  \item 只要该数为偶数，就先用 $2$ 去除；
  \item 然后依次试 $3$、$5$、$7$、\dots\ （只取质数）；
  \item 当商为 $1$ 时停止，并将各因子
        连同其指数一并写出。
\end{enumerate}
只需试除质数 $p$，且 $p^2$ 不超过当前数，
若这些质数都不能整除它，则它本身便是质数。
\end{method}

\begin{example}\label{ex:g9:arith:factorization}
对 $360$ 作质因数分解，每次做一次除法：
\[
360 = 2 \times 180, \quad
180 = 2 \times 90, \quad
90 = 2 \times 45, \quad
45 = 3 \times 15, \quad
15 = 3 \times 5,
\]
于是
\[
360 = 2 \times 2 \times 2 \times 3 \times 3 \times 5 = 2^3 \times 3^2
\times 5 .
\]
\end{example}

\begin{omfigure}
\begin{tikzpicture}[scale=0.78]
  \node (n360) at (0,0) {$360$};
  \node (n180) at (1.2,-1) {$180$};
  \node (n90) at (2.4,-2) {$90$};
  \node (n45) at (3.6,-3) {$45$};
  \node (n15) at (4.8,-4) {$15$};
  \node (n5) at (6.0,-5) {$5$};
  \node[omThm] (p2a) at (-1.2,-1) {$2$};
  \node[omThm] (p2b) at (0,-2) {$2$};
  \node[omThm] (p2c) at (1.2,-3) {$2$};
  \node[omThm] (p3a) at (2.4,-4) {$3$};
  \node[omThm] (p3b) at (3.6,-5) {$3$};
  \draw (n360) -- (n180); \draw (n360) -- (p2a);
  \draw (n180) -- (n90); \draw (n180) -- (p2b);
  \draw (n90) -- (n45); \draw (n90) -- (p2c);
  \draw (n45) -- (n15); \draw (n45) -- (p3a);
  \draw (n15) -- (n5); \draw (n15) -- (p3b);
\end{tikzpicture}

{\small $360$ 的因数树：每一步分出最小的质因数
（红色）。读取红色的叶子以及最后的 $5$：
$360 = 2^3 \times 3^2 \times 5$。}
\end{omfigure}

\begin{theorem}[欧几里得]\label{thm:g9:arith:euclid}
质数有无穷多个。
\end{theorem}

\begin{proof}
假设质数只有有限多个，设为 $p_1, p_2, \dots, p_k$，
并考虑
\[
N = p_1 \times p_2 \times \dots \times p_k + 1 .
\]
用 $N$ 除以任一 $p_i$，余数均为 $1$，因此没有任何 $p_i$ 整除 $N$。
但 $N \geq 2$ 至少有一个质因数
（见\cref{thm:g9:arith:factorization}） --- 一个不在该列表中的质数。
矛盾：任何有限列表都不可能包含全部质数。
\end{proof}

\section{最大公因数}

\begin{definition}[最大公因数]\label{def:g9:arith:gcd}
称能同时整除两个正整数的最大整数为\emph{最大公因数}\index{最大公因数}，
两个正整数 $a$、$b$ 的最大公因数记作 $\gcd(a, b)$。
当 $\gcd(a, b) = 1$ 时，称这两个整数为
\emph{互质}\index{互质}：它们除 $1$ 外不含其他公因数。
\end{definition}

\begin{example}
$18$ 的因数：$1, 2, 3, 6, 9, 18$。$24$ 的因数：$1, 2, 3, 4, 6,
8, 12, 24$。公因数为 $1, 2, 3, 6$；因此 $\gcd(18, 24) = 6$。
整数 $15$ 与 $28$ 互质。
\end{example}

\begin{proposition}[由质因数分解求最大公因数]\label{prop:g9:arith:gcdfact}
两个整数的最大公因数，等于同时出现在\emph{两个}分解式中的质数的乘积，
每个质数取其两个指数中的\emph{较小}者
作为指数。
\end{proposition}
\admitted

\begin{example}\label{ex:g9:arith:gcdfact}
$360 = 2^3 \times 3^2 \times 5$ 与 $84 = 2^2 \times 3 \times 7$。公共
质数：$2$（指数分别为 $3$ 与 $2$，保留 $2$）以及 $3$（指数分别为 $2$ 与
$1$，保留 $1$）。于是
\[
\gcd(360, 84) = 2^2 \times 3 = 12 .
\]
\end{example}

\begin{theorem}[欧几里得算法]\label{thm:g9:arith:euclidalgo}
若 $a = bq + r$ 是 $a$ 除以 $b$ 的带余除法（余数为 $r$），则
\[
\gcd(a, b) = \gcd(b, r).
\]
反复做除法直至余数为 $0$，则 $a$ 与 $b$ 的最大公因数是
\emph{最后一个非零余数}。
\end{theorem}

\begin{proof}
由 $a = bq + r$ 知，任何同时整除 $b$ 与 $r$ 的整数都整除 $bq + r = a$；
又由 $r = a - bq$ 知，任何同时整除 $a$ 与 $b$ 的整数都整除 $r$。于是
数对 $(a, b)$ 与 $(b, r)$ 恰有相同的公因数 ---
特别地，最大的公因数也相同。由于余数严格递减，
算法必然终止，且由 $\gcd(x, 0) = x$ 可得最后一个非零
余数。
\end{proof}

\begin{example}\label{ex:g9:arith:euclidalgo}
计算 $\gcd(1071, 462)$：
\begin{align*}
1071 &= 462 \times 2 + 147, \\
462 &= 147 \times 3 + 21, \\
147 &= 21 \times 7 + 0 .
\end{align*}
最后一个非零余数为 $21$，故 $\gcd(1071, 462) = 21$。
\end{example}

\begin{method}[彻底约分]\label{met:g9:arith:simplify}
要将 $\dfrac ab$ 化为既约分数：
\begin{enumerate}
  \item 计算 $d = \gcd(a, b)$，例如用欧几里得算法；
  \item 分子与分母同除以 $d$：
        $\dfrac ab = \dfrac{a \div d}{b \div d}$；
  \item 所得分数为\emph{既约分数}：其分子与
        分母互质。
\end{enumerate}
\end{method}

\begin{example}
$\dfrac{462}{1071} = \dfrac{462 \div 21}{1071 \div 21} = \dfrac{22}{51}$，
且 $\gcd(22, 51) = 1$，故为既约分数。
\end{example}

\section{练习}

\begin{exercise}[$\star$]\label{exo:g9:arith:1}
列出 $36$、$45$、$17$ 的全部因数。
\end{exercise}

\begin{exercise}[$\star$]\label{exo:g9:arith:2}
利用整除判别法，判断 $2\,346$ 是否能被
$2$、$3$、$4$、$5$、$9$ 整除。
\end{exercise}

\begin{exercise}[$\star$]\label{exo:g9:arith:3}
分别写出 $72$、$150$、$210$ 与 $121$ 的质因数分解。
\end{exercise}

\begin{exercise}[$\star$]\label{exo:g9:arith:4}
判断 $101$、$91$、$143$ 是否为质数，并利用
\cref{met:g9:arith:factorization}中的停止准则说明理由。
\end{exercise}

\begin{exercise}[$\star$]\label{exo:g9:arith:5}
用两种方法计算 $\gcd(48, 60)$：一是列出
公因数，二是由质因数分解求出。
\end{exercise}

\begin{exercise}[$\star\star$]\label{exo:g9:arith:6}
用欧几里得算法计算 $\gcd(255, 154)$，再计算
$\gcd(1053, 325)$。写出每一步的除法算式。
\end{exercise}

\begin{exercise}[$\star\star$]\label{exo:g9:arith:7}
将分数 $\dfrac{588}{504}$ 化为既约分数。（任选一种方法
求出最大公因数，再进行约分。）
\end{exercise}

\begin{exercise}[$\star\star$]\label{exo:g9:arith:8}
一位花店店主有 $84$ 朵玫瑰和 $126$ 朵郁金香，她想扎成
完全相同的花束，并将所有花都用上，且花束尽可能多。
她最多能扎成多少束花？每束花里各有几朵玫瑰和几朵郁金香？
\end{exercise}

\begin{exercise}[$\star\star$]\label{exo:g9:arith:9}
两班渡轮于 8:00 从同一码头发出。一班每 $24$ 分钟发出一次，
另一班每 $36$ 分钟发出一次。它们下一次同时发出的时刻是几点？
（求 $24$ 与 $36$ 的最小公倍数；
质因数分解会有帮助。）
\end{exercise}

\begin{exercise}[$\star\star\star$]\label{exo:g9:arith:10}
设 $n$ 为正整数。
\begin{enumerate}
  \item 证明 $\gcd(n, n+1) = 1$（相邻两个整数总是
        互质）。
  \item 由此推出分数 $\dfrac{n}{n+1}$ 总是既约分数。
\end{enumerate}
\end{exercise}

\section{问题：水罐、蝉与一百个储物柜}

\begin{problem}[{周末问题 --- 最大公因数决定两个水罐能量出哪些量；
质数保护蝉；保持开着的储物柜
编号恰为完全平方数}]
\label{pb:g9:arith:1}
三道看似谜语、实为算术的题目：
用无刻度的水罐量水（即伪装起来的最大公因数），昆虫
的生命周期演化为质数，以及一条著名的
一百个储物柜的走廊，其最终状态由
统计因数决定。全部内容都依赖于本章的工具：
整除性、质因数分解
（\cref{thm:g9:arith:factorization}）与欧几里得算法
（\cref{thm:g9:arith:euclidalgo}）。

\medskip
\noindent\textbf{第一部分 --- 水罐。}
你站在喷泉旁，手边有两个无刻度的水罐，容量分别为 $5$ L 和
$3$ L。允许的操作：将水罐注满，
将水罐完全倒空，把一个水罐的水倒入另一个，
直到源罐变空或目标罐装满。
\begin{enumerate}
  \item 恰好量出 $1$ L。（写出你的操作顺序，
        以及每次操作后两个水罐中的水量。）
  \item 恰好量出 $4$ L --- 一部著名动作片中的
        谜题。（可以用六步完成。）
  \item 从 $1$ 到 $8$ 升中，哪些整数升量可以量出
        （只用一个水罐，或分装在两个水罐中）？请补全
        这张清单，并利用已有的操作序列。
  \item 换两个水罐：$6$ L 和 $4$ L。试着量出 $1$ L --- 然后
        解释为何这是行不通的：检验三种允许的操作中，
        每个水罐中的水量始终是 $2$ 的倍数，
        因此任何能量出的水量都是偶数。
  \item 第 4 题的论证具有一般性：设两个水罐的容量分别为 $a$
        和 $b$ 升，则任何能量出的水量都是
        $\gcd(a, b)$ 的倍数。计算 $\gcd(6, 4)$ 与 $\gcd(5, 3)$，并
        说明该规律对每一对水罐分别预测了什么结果。
\end{enumerate}

\medskip
\noindent\textbf{第二部分 --- 喷泉旁的欧几里得。}
\begin{enumerate}[resume]
  \item 用欧几里得算法计算：$\gcd(91, 65)$ 与
        $\gcd(2\,026, 46)$。
  \item 用你自己的话解释，水罐中出现的水量为何本质上就是
        伪装起来的欧几里得余数：
        取容量分别为 $13$ L 和 $5$ L 的水罐，反复把小罐注满后倒入大罐
        （大罐装满时就将其倒空）。
        哪些新的水量最先出现？请将它们
        与欧几里得算法中的余数进行比较，
        以 $(13, 5)$ 为例。
  \item 推出冠军的答案：若水罐的容量为 $13$ 与 $5$ 升，
        你能否恰好量出 $1$ L？请用一行说明理由，
        结合第 5 题与 $\gcd(13, 5)$。
  \item 仿照
        \cref{exo:g9:arith:10}，简短证明 $\gcd(n, 2n + 1) = 1$
        对每个正整数 $n$ 成立。
        （$n$ 与 $2n + 1$ 的公因数必须整除什么？）
  \item 两辆公交车于 7:00 同时从起点站发车；一辆
        每 $12$ 分钟发一班，另一辆每 $18$ 分钟发一班。列出
        两辆车各自接下来的发车时刻，并求出它们
        再次同时发车的最早时刻。在这个例子上验证那条漂亮的规律：
        （最小公倍数）$\times$
        $\gcd$ $=$ 两数之积 --- 并
        再用 $5$ 和 $3$ 检验一次。
\end{enumerate}

\medskip
\noindent\textbf{第三部分 --- 蝉、因数与储物柜。}
\begin{enumerate}[resume]
  \item 北美的某些蝉只每 $17$ 年出土一次；
        假设某种捕食者的数量每 $4$ 年达到一次峰值。
        若两者今年同时发生，那么再过多少年，
        下一次出土恰与峰值重合？若蝉的周期改为 $16$ 年，问题相同 ---
        那么它们将多久遭到一次大规模捕食？
        请用一句话解释，为何进化会把周期推向
        \emph{质数}长度。
  \item 利用分解式 $360 = 2^3 \times 3^2 \times 5$，
        不必列出，直接数出 $360$ 的因数个数：一个因数对 $2$ 的指数
        有四种选择（$0, 1, 2, 3$），
        对 $3$ 与 $5$ 的指数各选一次，
        一共有多少个因数？
  \item 证明：在完全平方数
        $n = m^2$ 的质因数分解中，每个质数的指数都是\emph{偶数}。
        由此不计算任何平方根，即可推出 $360$ 不是
        完全平方数。
  \item 将 $d$ 与其伙伴配对，$n$ 的全部因数都参与配对。
        即 $\frac{n}{d}$（以 $n = 36$ 为例：$1 \leftrightarrow 36$，
        $2 \leftrightarrow 18$，$3 \leftrightarrow 12$，
        $4 \leftrightarrow 9$，$6 \leftrightarrow 6$）。何时一个因数恰与自己配对？
        由此推出判别准则：$n$ 有
        \emph{奇数}个因数，当且仅当 $n$ 是
        完全平方数。在 $36$ 与 $360$ 上检验这一结论。
  \item 一百个储物柜。编号 $1$ 至 $100$ 的储物柜开始时均为关闭状态。
        学生 $1$ 切换每一个储物柜的状态；学生 $2$ 切换
        $2, 4, 6, \dots$ 号储物柜；学生 $k$ 切换
        $k$ 的倍数号储物柜；依此类推，直到学生 $100$。
        请解释哪些学生会切换 $n$ 号储物柜，它被
        切换了多少次，并利用第 14 题 --- 精确确定
        哪些储物柜最终是打开的。共有多少个打开的储物柜？
\end{enumerate}
\end{problem}
